\subsection{Semi-norms in quotient spaces and in product spaces}

Let E be a topological vector space over K, whose topology is defined by \(\Gamma\) , a set of semi-norms. Clearly, the restrictions of the semi-norms of \(\Gamma\) to a vector sub-space M of E, define the topology induced on M by that of E.

Let \(\phi\) be the canonical mapping of E on the vector quotient space E/M. We show that, for every semi-norm \(p\) on E, the function

\[
\dot {p} (z) = \inf _ {\phi (x) = z} p (x)\tag{3}
\]

is a semi-norm on E/M. In fact, it is clear that \(\dot{p}\) satisfies the condition \((\mathrm{SN}_{\mathrm{I}})\) ; on the other hand, if \(z'\) , \(z''\) are two vectors of E/M, we have :

\[
\begin{array}{r l} \inf _ {\phi (x) = z ^ {\prime} + z ^ {\prime \prime}} p (x) & \leqslant \inf _ {\phi (x ^ {\prime}) = z ^ {\prime}, \phi (x ^ {\prime \prime}) = z ^ {\prime \prime}} p (x ^ {\prime} + x ^ {\prime \prime}) \\ & \leqslant \inf _ {\phi (x ^ {\prime}) = z ^ {\prime}, \phi (x ^ {\prime \prime}) = z ^ {\prime \prime}} \left(p (x ^ {\prime}) + p (x ^ {\prime \prime})\right) \\ & = \inf _ {\phi (x ^ {\prime}) = z ^ {\prime}} p (x ^ {\prime}) + \inf _ {\phi (x ^ {\prime \prime}) = z ^ {\prime \prime}} p (x ^ {\prime \prime}) \end{array}
\]

which shows that \(\dot{p}\) verifies \((\mathrm{SN}_{\mathrm{II}})\) . We say that \(\dot{p}\) is the quotient semi-norm of p by M.

The same reasoning proves that, if p is an ultra-semi-norm, then so also is \(\dot{p}\) .

This being so, we have (in the notation of No.~2)

\[
\phi (\mathrm{W} (p, \alpha)) = \mathrm{W} (\dot {p}, \alpha).\tag{4}
\]

for every \(\alpha > 0\) . In fact, to say that \(\dot{p}(z) < \alpha\) , means that there exists \(x \in \mathbf{E}\) such that \(\phi(x) = z\) and \(p(x) < \alpha\) , from which the relation (4) follows.

We deduce from this, that, if the set of semi-norms \(\Gamma\) is directed(II, p.~3, Remark 3), then the quotient topology on E/M is defined by the set of semi-norms \(\dot{p}\) , when p varies in \(\Gamma\) .

If N is the closure of \(\{0\}\) in E, the topology of E/N is defined by the quotient seminorms \(\dot{p}\) , where p varies in \(\Gamma\) (even if \(\Gamma\) is not filtered): here \(\dot{p}(\dot{x}) = p(x)\) for every x belonging to the class \(\dot{x} \mod N\) . Note that E/N is none other than the Hausdorff space associated with E (I, p.~4).

Let \(\mathbf{E}\) be a vector space over \(\mathbf{K}\) and \((\mathrm{E}_{\iota})_{\iota \in \mathbf{I}}\) be a family of vector spaces over \(\mathbf{K}\) , where \(\mathrm{E}_{\iota}\) has the topology \(\mathcal{T}_{\iota}\) defined by a set of semi-norms \(\Gamma_{\iota}\) . For each \(\iota \in \mathbf{I}\) , let \(f_{\iota}\) be a linear mapping of \(\mathbf{E}\) in \(\mathrm{E}_{\iota}\) ; clearly when \(p_{\iota}\) varies in the set \(\Gamma_{\iota}\) , then the \(p_{\iota} \circ f_{\iota}\) form a set \(\Gamma_{\iota}'\) of semi-norms on \(\mathbf{E}\) . The topology \(\mathcal{T}\) on \(\mathbf{E}\) , defined as being the coarsest of all those which make all the mappings \(f_{\iota}\) continuous (I, p.~9) is then defined by the set of semi-norms \(\Gamma' = \bigcup_{\iota \in \mathbf{I}} \Gamma_{\iota}'\) , this follows from the definition of neighbourhoods of 0 for \(\mathcal{T}\) (GT, I, § 2.3, prop. 4).

If the \(p_{\iota}\) are ultra-semi-norms, then so are the \(p_{\iota} \circ f_{\iota}\) .

Let E be a vector space over K, with the topology T defined by a family of semi-norms \((p_{\iota})_{\iota \in I}\) ; for every \(\iota \in I\) , let \(T_{\iota}\) be the topology defined by the single semi-norm \(p_{\iota}\) , and denote by \(E_{\iota}\) the space obtained from E using the topology \(T_{\iota}\) . Then the topology T is the inverse image by the diagonal mapping \(\Delta : E \to \prod_{\iota \in I} E_{\iota}\) of the product topology on \(\prod_{\iota \in I} \mathrm{E}_{\iota}\) (I, p.~9, prop. 7). For each \(\iota \in I\) , write \(\mathbf {N} _ {\iota}\) for the closure of \(\{0 \}\) in \(\mathrm{E} _ {\iota}\) , and by \(F_{\iota} = E_{\iota}/N_{\iota}\) , the normed space defined by the norm \(\hat{p}_{\iota}\) corresponding to \(p_{\iota}\) (II, p.~4, formula (3)); if \(\phi_{\iota}: E_{\iota} \to F_{\iota}\) is the canonical mapping and \(\phi : (x_{\iota}) \mapsto (\phi_{\iota}(x_{\iota}))\) the product mapping, we know that the product topology on \(\prod_{i \in I} E_{i}\) is the inverse image by \(\phi\) of the product topology on \(\prod_{\iota \in I}F_{\iota}\) (GT, II, § 3.9, prop. 18). The topology \(\mathcal{T}\) is, therefore, the inverse image under the composite mapping \(\phi \circ \Delta\) of the product topology on \(\prod_{\mathfrak{t}\in \mathbf{I}}\mathrm{F}_{\mathfrak{t}}\) . In particular, if \(\mathcal{T}\) is Hausdorff then it follows from II, p.~3, prop. 2 that \(\Delta^{-1}(\prod_{\iota \in I} \mathbf{N}_{\iota}) = \{0 \}\) , therefore that the mapping \(\phi \circ \Delta\) is injective, therefore :

PROPOSITION 3. --- Every Hausdorff topological vector space E over K, whose topology is defined by a set of semi-norms, is isomorphic to a sub-space of a product of Banach spaces.

If, further, the topology of E is defined by an enumerable set of semi-norms, then E is metrisable (I, p.~16).

